Great question! Converting an **arbitrary state \(|\psi_1\rangle\) to another arbitrary state \(|\psi_2\rangle\)** is a classic scenario, and it’s a bit deeper than just state preparation from \(|0\rangle^{\otimes n}\). Here are clear answers and practical strategies:

---

# 1. **Any two pure states can be connected by a unitary**

Suppose you have:
- \(|\psi_1\rangle\) (your current state)
- \(|\psi_2\rangle\) (your target state)

**There always exists a unitary \(U\) such that**
\[
U|\psi_1\rangle = |\psi_2\rangle
\]

This is a direct result from linear algebra (they are both normalized vectors in Hilbert space), though \(U\) is generally *not unique*.

---

# 2. **Constructing the required unitary**

Suppose you can implement an arbitrary unitary operation.

### The Householder reflection approach

1. **Define the Householder operator:**
   \[
   U = I - 2|\phi\rangle\langle \phi|
   \]
   where
   \[
   |\phi\rangle = \frac{|\psi_1\rangle - |\psi_2\rangle}{\| |\psi_1\rangle - |\psi_2\rangle \| }
   \]
   This is a reflection that maps \(|\psi_1\rangle\) to \(|\psi_2\rangle\): \(U|\psi_1\rangle = |\psi_2\rangle\) (up to a phase).

   - For details, see [Nielsen & Chuang, Eqn 4.28](https://books.google.com/books?id=65FqEKQOfP8C&pg=PA186&lpg=PA186).

2. **Building the circuit**
   - In principle, you could decompose this Householder unitary into elementary gates, using known decomposition techniques (like those described above).

---

### Or: **State “un-preparation” + preparation**

Most practical toolkits (like Qiskit) provide methods to go **from \(|0…0⟩\) to arbitrary target states**.

If you have a circuit (think: a sequence of gates) that sends \(|0\rangle\) → \(|\psi_1\rangle\), call it \(U_1\):

\[
U_1|0\rangle = |\psi_1\rangle
\]

Similarly, let \(U_2\) prepare \(|\psi_2\rangle\):

\[
U_2|0\rangle = |\psi_2\rangle
\]

Then the operation:

\[
U = U_2 U_1^\dagger
\]

maps
\[
U |\psi_1\rangle = U_2 U_1^\dagger |\psi_1\rangle = U_2 |0\rangle = |\psi_2\rangle
\]

**So, if you have circuits for state-preparing \(|\psi_1\rangle\) and \(|\psi_2\rangle\), implement \(U_1^\dagger\) followed by \(U_2\):**
- 1. Undo \(|\psi_1\rangle\) back to \(|0\rangle\) via \(U_1^\dagger\)
- 2. Prepare \(|\psi_2\rangle\) via \(U_2\)

---

# 3. **How to do it in practice (with Qiskit)**

Let’s say you have **state vectors** for \(|\psi_1\rangle\) and \(|\psi_2\rangle\).

You can use Qiskit’s **Initialize** gate for \(|\psi_1\rangle\) and \(|\psi_2\rangle\):

```python
from qiskit import QuantumCircuit
from qiskit.circuit.library import Initialize
import numpy as np

# psi_1 and psi_2 are *normalized* numpy arrays of 2**n amplitudes
psi_1 = np.array([...])
psi_2 = np.array([...])

n = int(np.log2(len(psi_1)))

qc = QuantumCircuit(n)

# Un-preparation: psi_1^dagger (adjoint of the gate that prepares psi_1 from |0>)
init1 = Initialize(psi_1)
init1_circ = init1.gates_to_uncompute()   # The inverse circuit

# Preparation: prepare psi_2 from |0>
init2 = Initialize(psi_2)

# Compose: first uncompute psi_1, then compute psi_2
qc.append(init1_circ, range(n))
qc.append(init2, range(n))
```

- If the qubits **start in \(|\psi_1\rangle\)**, this circuit will transform them to \(|\psi_2\rangle\).
- \(Initialize\) decomposes to basic gates and is fully supported for this kind of task.

---

# 4. **Summary Table**

| Step                                          | Qiskit Method                                      |
|-----------------------------------------------|----------------------------------------------------|
| Prepare \(|\psi_1\rangle\) from \(|0\rangle\) | `Initialize(psi_1)`                                |
| Uncompute \(|\psi_1\rangle\)                  | `Initialize(psi_1).gates_to_uncompute()`           |
| Prepare \(|\psi_2\rangle\) from \(|0\rangle\) | `Initialize(psi_2)`                                |
| Convert \(|\psi_1\rangle\) → \(|\psi_2\rangle\) | `init1_adjoint` + `init2` as above                 |

---

## **In summary:**
1. **Any pure state can be mapped to any other pure state by a unitary.**
2. **You can build a gate sequence for this by “unpreparing” the initial state to \(|0...0⟩\), then preparing the target.**
3. **In Qiskit, use the inverse of `Initialize(psi_1)` to undo, then `Initialize(psi_2)` to redo.**

---

**Want a concrete numerical example? Just give explicit \(|\psi_1⟩\), \(|\psi_2⟩\)!**